\chapter{Expressions and control flow}\label{ch-control}

The programs of the previous chapter ran from top to bottom. This chapter
introduces the constructs that choose between alternatives and repeat
actions, and the rules that say where a declared name can be used. It
begins with the idea that makes Algol 68 different from most languages of
its time and from many of ours: there is no separate category of
statements. Every construct, including a block, a choice and a loop, is an
expression.

\section{Everything is an expression}\label{sec-everything}

In C, Java or Python, a program is made of statements, and expressions
are one kind of statement. Algol 68 has only expressions, which the Report
calls \rr{units}. Constructs that look like statements, such as an
assignment or a call of \code|print|, are expressions too; most of them
simply have a value that nobody uses.

A block, written either \code|begin|\ldots\code|end| or
\code|(|\ldots\code|)|, contains a sequence of declarations and expressions
separated by semicolons, and its value is the value of the last
expression. A block can therefore stand wherever a value is needed:

\example{ch03/block-value}

The two spellings of a block mean exactly the same; \code|begin| and
\code|end| are usually preferred for long blocks, parentheses for short
ones. The Report calls a block a \rr{closed clause} and the sequence inside
it a \rr{serial clause}. Rust and the functional languages work the same
way; in C, the closest equivalent is the ``statement expression'' of the
GNU C compiler.

A block must end with an expression, not with a declaration, since it
must have a value. The block \code|begin int x = 1 end| is an error
(``clause does not yield a value''), even in a place where its value is
not used.

The value of a block is computed before the block ends, so a block may
yield the value of one of its own variables: \code|int n = (int r; read (r); r)|
reads a number through a local variable and yields it. What a block may
not yield is the variable itself, a reference to a cell that disappears
with the block; Chapter~\ref{ch-advanced} explains this rule about
lifetimes.

\section{Sequences and discarded values}\label{sec-voiding}

When an expression is followed by a semicolon, its value is discarded. The
Report calls this \rr{voiding}, and it is one of the implicit conversions:
the value is converted to the type \code|void|, which has a single,
uninteresting value. Discarding is what happens to the value of every
assignment and every call of \code|print| in the programs so far:

\example{ch03/voiding}

The value of \code|x := x + 1| is the reference \code|x|
(Section~\ref{sec-assignment}); it is discarded, and so is the value of
\code|x + 1| in the line before, which therefore does nothing useful.

\section{Scope}\label{sec-scope}

A name declared in a block can be used anywhere in that block, including
in blocks nested inside it, but not outside. The part of the program in
which a declaration is valid is its \term{scope} (the Report says
\rr{range} for the block, and \rr{reach} for the block without the
blocks nested in it). A nested block can declare a
name that is already declared outside; inside the nested block, the new
declaration hides the outer one:

\example{ch03/scope}

The same name cannot be declared twice in the same block; this is an
error found by the compiler (``multiple declaration of identifier'').

Strictly speaking, the scope of a declaration is the whole block,
including the part before the declaration. The name can be mentioned
there, but its value cannot be used, since the declaration has not been
carried out yet. A program that uses a name before its declaration is
therefore wrong, but the compilers do not detect it. What happens at run
time depends on the value. An \code|int| constant used before its
declaration is 0 in ga68 and a68j; Algol 68 Genie stops the program with
a message. A function that has no value yet, such as a variable for
functions (Chapter~\ref{ch-procs}) before its first assignment, makes the
ga68 program crash; a68j and Algol 68 Genie stop with a message, a68j with
``attempt to use an undefined value''. A function declared with a
function literal, the usual form of Chapter~\ref{ch-procs}, is an
exception: in all three it exists from the start of the block and can be
called even before its declaration. The rule matters for functions that
call each other (Chapter~\ref{ch-procs}); for everything else, a
declaration should simply come before its uses.

A word about the term: in this book, \emph{scope} means the region of the
program text in which a name can be used, as in every current language.
The Report uses the word \rr{scope} for something else, the lifetime of a
value, which is the subject of Chapter~\ref{ch-advanced}.

\section{Unordered evaluation}\label{sec-collateral}

A list of expressions separated by commas and enclosed in parentheses is
evaluated in an unspecified order:

\example{ch03/collateral}

Here the two assignments may be carried out in either order, or even
interleaved; the program must not depend on the order. The Report calls
this \rr{collateral elaboration}, and the list a \rr{collateral clause}.
The same holds wherever several expressions separated by commas are all
evaluated: the operands of an operator, the arguments of a procedure
call, the items of a list given to \code|print|, the elements of an array
literal, and the declarations in \code|int a = f, b = g|. (The commas
between the alternatives of a \code|case| are different: only one
alternative is evaluated.) A program
whose result depends on the order in which \code|f| and \code|g| are
called is wrong, even if it happens to work with a particular compiler.
Some places without commas have the same freedom: the two sides of an
assignment, an array and its indexes in \code|a[i]|, and a function and
its arguments in a call.

A semicolon, by contrast, always means ``first this, then that''.

\begin{pitfall}
Code that looks harmless can depend on the order. In
\code|(k +:= 1) * k|, the right operand \code|k| may be read before or
after the left one has incremented it. In \code|int a = 1, b = a + 1|,
the two declarations are carried out in an unspecified order, so
\code|b| may be computed before \code|a| exists; C programmers write this
all the time, and C defines it, but Algol 68 does not. The same holds for
two reads in one expression, as in \code|read_int - read_int| with
Algol 68 Genie's \code|read_int|, and for
\code|a[i] := (i +:= 1)| with \code|i| = 1, which may store 2 into
\code|a[1]| or into \code|a[2]|. ga68 and a68j evaluate from left to
right (here into \code|a[1]|), so testing does not reveal the error. Writing \code|int a = 1; int b = a + 1| with
a semicolon makes the order definite. The items of a \code|print| list
are evaluated in an unspecified order, but they are always printed in
the order in which they are written.
\end{pitfall}

\section{\textbf{if} expressions}\label{sec-if}

The conditional expression has the form
\begin{lstlisting}
if condition then expressions else expressions fi
\end{lstlisting}
where \code|fi|, \code|if| spelled backwards, closes the construct. The
\code|else| part may be omitted, and further conditions may be tested
with \code|elif|. Like a block, an \code|if| has a value, so it serves
both as a statement and as the conditional expression \texttt{c ? a : b}
of C:

\example{ch03/if}

The last line shows the brief form: \code{(c | a | b)} means
\code|if c then a else b fi|, and \code{|:} stands for \code|elif|. The
brief form is convenient in short expressions; for anything longer, the
words are easier to read.

The condition is itself a sequence and may contain declarations, which are
then visible in all parts of the \code|if|:
\begin{lstlisting}
if int k = n * n; k > 100 then print (k) fi
\end{lstlisting}

If the \code|else| part is omitted and the condition is false, the
\code|if| does nothing, which is what is wanted in a statement. If such
an \code|if| is used for its value, the value is undefined. The Report
calls this undefined value \code|skip| (Section~\ref{sec-skip}).

When the branches have different types, the compiler looks for a common
type to which all branches can be converted, and converts them. If
one branch is an \code|int| and the other a \code|real|, the common type is
\code|real|:

\example{ch03/balancing}

The Report calls this \rr{balancing}. It is what C does for
\texttt{c ? 1 : 2.5}, and what the type inference of a modern language
does when it finds the least common type of the branches.

Since an \code|if| yields a value of any type, it can also yield a
variable, and stand on the left-hand side of an assignment:

\example{ch03/if-target}

\noindent sets \code|b| to 0. Here both branches yield a variable of
type \code|ref int|. With a \code|real| variable \code|x| and an
\code|int| variable \code|n|, \code|(c | x | n) := 4| is rejected: the
branches yield a \code|ref real| and a \code|ref int|, and neither can be
converted into the other.

\section{\textbf{case} expressions}\label{sec-case}

The \code|case| expression selects one of several alternatives by an
integer: the first alternative for 1, the second for 2, and so on. If the
integer is out of range, the \code|out| part is chosen; without an
\code|out| part, nothing happens (or the value is undefined).

\example{ch03/case}

The alternatives are separated by commas, and each can be a single
expression or a block. The Report requires at least two alternatives (for
one, an \code|if| does the job); the three implementations also accept
one. The brief form is \code{(i | a, b, c | d)}. A second
selection can follow the \code|out| with \code|ouse|, much as \code|elif|
follows the \code|then| part of an \code|if|.

This \code|case| corresponds to the \texttt{switch} of C, except that the
alternatives are numbered from 1 and cannot be labeled with arbitrary
constants. Chapter~\ref{ch-advanced} introduces a second, more powerful form of
\code|case| that selects by the type of a value.

\section{Loops}\label{sec-loops}

Algol 68 has a single loop construct with several optional parts:
\begin{lstlisting}
for i from first by step to last while condition do expressions od
\end{lstlisting}
Every part except \code|do|\ldots\code|od| may be left out. The parts
have the following meaning:
\begin{itemize}
\item \code|for i| names the \term{loop variable}; without it, the loop
  still counts, but the count has no name.
\item \code|from first| gives the initial value, 1 if omitted.
\item \code|by step| gives the increment, 1 if omitted; it may be
  negative.
\item \code|to last| gives the final value; without it, there is no limit.
\item \code|while condition| is tested before each iteration; the loop
  ends when it is false.
\end{itemize}
The values of \code|first|, \code|step| and \code|last| are computed once,
before the first iteration. The loop variable is an \code|int| constant,
not a variable: it is declared anew for each iteration, it cannot be
assigned to, and it does not exist after the loop. The usual forms of the
loop result from the parts that are present:

\example{ch03/for-loops}

\code|to 3 do| \ldots \code|od| repeats the body three times. A loop with
only a \code|while| part is the \texttt{while} loop of other languages;
Euclid's algorithm for the greatest common divisor is a classic example:

\example{ch03/while-loop}

The \code|for| and \code|while| parts can be combined: the loop then ends
when the counter passes the final value or the condition becomes false,
whichever comes first. Like the condition of an \code|if|, the
\code|while| part may contain declarations, which are visible in the
body. This replaces the \texttt{do}\ldots\texttt{while} loop of C, whose
body runs before its test: the part of the body that must run first is
written in the \code|while| part.

\example{ch03/for-while}

Finally, \code|do|\ldots\code|od| alone repeats forever; the loop is left
with a jump (Section~\ref{sec-jumps}). A loop has no value: its type is
\code|void|.

Since the loop variable is a constant, an attempt to assign to it fails
with the message we met in Chapter~\ref{ch-values}:

\example{ch03/loop-variable}

A few details of the loop are worth knowing. Since the final value is
computed once, changing a variable that appears in the \code|to| part
inside the body does not change the number of iterations. The
\code|while| part is not evaluated any more once the counter has passed
the final value, which matters if it reads input or has other effects.
With a step of 0 the final value is ignored, and the loop runs forever
unless the \code|while| part stops it. (ga68 and a68j instead skip such
a loop if the start value is above the final value.)

The loop variable is always an \code|int|; there is no loop over real
numbers with a step of 0.1. That is just as well, since real numbers are
approximations, and adding 0.1 ten times does not give exactly 1.0:

\program{ch03/real-sum}

\noindent prints
\outputof{ch03/real-sum}

The right way is to count with an integer and compute each real value
from the counter, as in the second loop, so that the errors do not
accumulate. For the same reason, real numbers should not be compared with
\code|=|, but tested for a small difference, as in
\code|ABS (sum - 1.0) < 1e-6|. With the 32-bit \code|real| of ga68 and
a68j, which has about 7 decimal digits, the problem appears sooner than
with the 64-bit \code|double| that most programmers are used to.

The \code|while| part is the place for an iteration that must compute a
value before it can decide whether to continue. Newton's method for the
square root improves an approximation until the improvement is small
relative to the value:

\example{ch03/newton}

\noindent prints \code|1.41421 1000.00 0.50000|. The test is relative,
\code|1e-6 * x|, because a fixed tolerance such as \code|1e-6| is far too
coarse for very small roots, and for large ones, with only 7
significant digits, it is reached only if the iteration happens to
arrive at an exact fixed point.

\begin{implnote}
Algol 68 Genie has two extensions that appear in Rosetta Code programs:
\code|downto| counts downwards (\texttt{FOR i FROM 3 DOWNTO 1}), and an
\texttt{UNTIL} part before \code|od| ends the loop after the body if its
condition is true. ga68 and a68j have neither; \code|by -1| and a
\code|while| part do the same.
\end{implnote}

\section{\textbf{skip}}\label{sec-skip}

The reserved word \code|skip| stands for an undefined value of whatever
type is needed. It is used where the syntax requires an expression but
no particular value is meaningful, for example as an alternative of a
\code|case| that has nothing to do, or as the body of a loop that does
all its work in the \code|while| part:

\example{ch03/skip-while}

\noindent The loop finds the smallest \code|n| whose square is at least
1000, and prints 32. The semicolon in the \code|while| part makes sure
that \code|n| is incremented before it is squared; the shorter
\code|(n +:= 1) * n < 1000| would depend on the order of the operands.

\section{Jumps}\label{sec-jumps}

An expression in a sequence can be given a \term{label}, a name followed
by a colon, and the \code|goto| expression continues the program at the
labeled expression. \code|go to| may also be written as two words, and in
many places the label alone is enough. A typical use is to leave a loop
that runs forever:

\example{ch03/jumps}

A jump can only go to a label in the same block or in an enclosing
block, never into a block, a loop or a branch, since the label is not
visible there (Section~\ref{sec-scope}).

A jump never yields a value, so it is acceptable wherever the type is
determined by the context or by the other branches of an \code|if| or a
\code|case|. This makes it easy to give up when the data are bad:

\example{ch03/jump-value}

\noindent prints \code|2.00|, \code|1.50| and \code|negative input|. The
\code|if| has the type \code|real|, and its \code|else| part, a jump, is
acceptable as a \code|real|. The predefined label \code|stop| ends the
whole program.

One case is different. Where a function without parameters is expected,
as for a parameter of type \code|proc void|, a jump is not carried out:
it becomes a function that carries out the jump when it is called. A
function can thus be given an exit to take when something goes wrong:

\example{ch03/error-exit}

\noindent The first call returns normally. In the second, \code|mean|
calls \code|on_empty|, and the jump leaves \code|mean| and continues at
the label \code|no_data|. In the same way,
\code|[] proc void switch = (l1, l2, l3)| is an array of functions that jump, and
\code|switch[i]| jumps to the \code|i|-th label, as the switch of Algol
60 did.

Labels have two restrictions. In a block, no declaration may follow a
label, so that a jump can never skip a declaration whose name is used
after the label; a label before a declaration is rejected with a message
about an invalid sequence. And the condition of an \code|if|, a
\code|case| or a \code|while| may not contain labels.

A block may also end before its last expression. The reserved word
\code|exit| after an expression ends the sequence with that expression's
value; the part after \code|exit| can only be reached by a jump to a label
that follows it:

\example{ch03/exit}

Jumps were common in programs of the 1960s, and Algol 68 still has them.
With its general loop and the value-yielding \code|if| and \code|case|
they are rarely needed, and the programs in this book use them only to
leave a loop early. The keyword \code|exit| should not be confused with
the \texttt{exit} or \texttt{break} of other languages: it does not leave
a loop.

\section{The structure of a program}\label{sec-program}

A program is a block, or more generally any construct in brackets: an
\code|if|, a \code|case| or a loop is a program too. Everything it uses
that it does not declare itself
comes from the standard library, which can be thought of as a block
enclosing the program. In this book, a program is preceded by
\code|access Genie_Transput|, which places the declarations of that module
around the program as well.

\begin{implnote}
Algol 68 Genie also accepts a program without the enclosing
\code|begin| and \code|end|, a convenience that many Rosetta Code
programs use. ga68 and a68j require the brackets.
\end{implnote}

\section{A complete program}\label{sec-control-complete}

The program below lists the prime numbers below 60. For each number
\code|n|, it tries the divisors from 2 upwards while no divisor has been
found and the divisor squared does not exceed \code|n|. The short-circuit
\code|andth| (Section~\ref{sec-logic}) is not essential here, but it
expresses the intent: once \code|prime| is false, the rest of the
condition does not matter.

\example{ch03/primes}

\section*{Exercises}
\addcontentsline{toc}{section}{Exercises}

The \hyperref[ans-ch3]{answers} are at the end of the book.

\begin{enumerate}
\item What are the values of \code|(int a = 3; a * a + 1)| and of
  \code{(2 > 1 | 10 | 20) + 1}?
\item What does the following program print?
\begin{lstlisting}
int x = 1;
begin int y = x + 1; begin int x = y * 10; print (x) end; print (x) end
\end{lstlisting}
\item How many times is the body of each loop executed?\\
  (a) \code|for i from 5 to 1 do skip od|\\
  (b) \code|for i from 1 by 2 to 10 do skip od|\\
  (c) \code|for i from 10 by -2 to 1 do skip od|\\
  (d) \code|to 0 do skip od|
\item Write \code{(x > 0 | "plus" |: x < 0 | "minus" | "zero")} with
  reserved words instead of the brief form.
\item What type has \code{(n > 0 | 1 | 0.5)}, and why?
\item Why may a program not rely on the order of the two calls in
  \code|f + g|, if \code|f| and \code|g| print something?
\item Write a loop that prints the powers of 2 up to 1000 without using a
  \code|for| part.
\item What does \code|int k := 3; to k do k *:= 2; print (k) od| print?
\item What does \code|for i from 1 by 0 to 3 do print (i) od| do?
\item Is \code|string name = case k in "one", "two", "three" out goto bad esac|
  a correct declaration, if \code|bad| is a label of the enclosing serial
  clause? What happens for \code|k = 5|?
\item With \code|real x := 1.5; int n := 2; bool b = x > n|, the declaration
  \code{real v = (b | x | n)} is accepted, but
  \code{ref real r = (b | x | n)} is rejected. Why?
\item What is wrong with \code|int a = 1, b = a + 1|, and how is it
  corrected?
\item The loop \code|int n; while read (n); n < 1 AND n > 9 do ... od| is
  meant to repeat until a number from 1 to 9 is read. Why does it stop
  after the first number, whatever it is?
\end{enumerate}

\section*{Programming problems}
\addcontentsline{toc}{section}{Programming problems}

\begin{enumerate}
\item Print a multiplication table for the numbers 1 to 9, with the
  numbers aligned in columns.
\item Read numbers until a 0 is read, and print their count, their sum and
  their largest value.
\item The Collatz sequence starting at $n$ continues with $n/2$ if $n$ is
  even and with $3n+1$ if $n$ is odd, until it reaches 1. Print the
  starting value below 1000 with the longest sequence.
\item Print the perfect numbers below 10000 (numbers that are equal to
  the sum of their divisors other than themselves, such as
  $6 = 1 + 2 + 3$).
\item Print a calendar for a month, given the number of days and the
  weekday of the first day.
\item \textbf{Number bases.} Read a number written in a base between 2
  and 16 (the base first, then the digits, with \texttt{A} to
  \texttt{F} for 10 to 15) character by character, and print its value.
  Report a digit that is not valid for the base and stop with
  \code|goto stop|.
\item \textbf{Brackets.} Read a line and decide whether its round
  brackets are properly nested. Why can the same method not check two
  kinds of brackets, \texttt{()} and \texttt{[]}, without an array?
\item \textbf{Arithmetic without brackets.} Read an expression of
  integers with \texttt{+}, \texttt{-} and \texttt{*}, such as
  \texttt{2+3*4-5}, and evaluate it with the usual precedence. Keep a
  running sum and a pending product, without arrays.
\item \textbf{Fibonacci and overflow.} Print the Fibonacci numbers up
  to the largest one that fits in an \code|int|, testing before each
  addition whether it would overflow (\code|a > max_int - b|). How many
  are there with a 32-bit and with a 64-bit \code|int|?
\item \textbf{Cube roots.} Compute cube roots with Newton's iteration
  $x_{n+1} = (2 x_n + a / x_n^2) / 3$, with a relative test for the end,
  and compare the results with \code|exp (ln (a) / 3)|. What happens for
  negative \code|a|?
\item \textbf{Prime factors.} Print the factorization of a number as a
  product of prime powers, such as \texttt{360 = 2\^{}3 * 3\^{}2 * 5}.
\item \textbf{Armstrong numbers.} 153 equals the sum of the cubes of its
  digits. Find all three-digit numbers with this property.
\end{enumerate}
